% related_work.tex

\section{Related Work}

\paragraph{Multi-turn jailbreak attacks.}
All multi-turn jailbreak families share a persuasion substrate~\citep{zeng2024pap,hagendorff2026reasoning,li2024mhj}, yet each operationalizes it differently: Crescendo~\citep{russinovich2024crescendo} gradually escalates, FITD~\citep{fitd2024} exploits commitment-consistency bias, DD~\citep{deceptivedelight2024} combines emotional manipulation with role assignment, ActorAttack~\citep{ren2024actorattack} steers through multiple personas, and automated methods (PAIR~\citep{chao2023pair}, TAP~\citep{mehrotra2024tap}) iteratively refine adversarial prompts, extending the automated red-teaming paradigm~\citep{perez2022redteaming}. \citet{wei2023jailbroken} provide a unifying taxonomy; RED QUEEN~\citep{jiang2024redqueen} constructs concealed multi-turn scenarios across 14 harmful categories; benchmarks~\citep{mazeika2024harmbench,chao2024jailbreakbench} standardize single-turn comparison, while MultiBreak~\citep{song2026multibreak} targets multi-turn evaluation. This shared substrate with diverse surface strategies defines our evaluation: a detector must generalize across families, not merely succeed on its training distribution.

\paragraph{Multi-turn defenses and production safety.}
No existing defense evaluates cross-attack transfer. Recent conversation-level systems pair fine-tuned BERT~\citep{devlin2019bert} encoders with sequence-level aggregation via GRU hierarchies (DeepContext~\citep{deepcontext2026}), turn-intention inference (BIID~\citep{tong2025biid}), graph neural networks (G-Guard~\citep{gguard2025}), or predictive generation (SafeDream~\citep{safedream2026}), yet each evaluates on a single attack family. Representation-level approaches include X-Boundary~\citep{lu2025xboundary}, JBShield~\citep{zhang2025jbshield} via activated concept analysis, gradient-space analysis (GradSafe~\citep{xie2024gradsafe}), circuit breakers via representation engineering~\citep{zou2024circuitbreakers}, latent activation probing~\citep{kulkarni2026latentadversarial,kadali2026trace}, and vision-language detection~\citep{hua2026rcs}; these require white-box access and can be brittle under adaptive attacks~\citep{derya2026revisitingjbshield}. Our work uses black-box classification with systematic cross-attack OOD evaluation. Per-turn classifiers (Llama Guard~\citep{inan2023llamaguard,metallamaguard3_2024}, ShieldGemma~\citep{zeng2024shieldgemma}, WildGuard~\citep{han2024wildguard}, Constitutional Classifiers~\citep{anthropic2025constitutional}) cannot model multi-turn trajectories, and perturbation-based defenses~\citep{robey2024smoothllm} and safety-aware decoding~\citep{xu2024safedecoding} target single-turn token-level attacks. Our ablation fills this gap: decomposition across thirteen configurations under OOD evaluation identifies domain exposure as the dominant factor behind cross-attack transfer.

\paragraph{Computational persuasion and discourse.}
Persuasion and discourse frameworks offer interpretable vocabulary for attack strategies, but whether such structure aids detection has not been empirically disentangled. Foundational work (RST~\citep{mann1988rst}, PDTB~\citep{prasad2008pdtb}, dialogue acts~\citep{stolcke2000dialogueact}, Speech Act Theory~\citep{austin1962speechacts,searle1969speechacts}) models pragmatic structure at different granularities. Building on Cialdini's principles~\citep{cialdini2001influence}, PAP~\citep{zeng2024pap} identifies 40 persuasion techniques across jailbreak families, and \citet{hagendorff2026reasoning} find identifiable strategies in 97\% of successful attacks; propaganda analysis~\citep{dasanmartino2019fine,dasanmartino2020semeval} and persuasive dialogue~\citep{wang2019persuasion} apply similar technique-level classification. Our ablation disentangles these roles: the taxonomy enables interpretable strategy characterization, but detection operates through domain-adapted embeddings regardless of label correctness.

\paragraph{Domain-adaptive pretraining.}
Continued pretraining on domain text yields consistent gains~\citep{gururangan2020dontstop}, as shown in biomedicine~\citep{lee2020biobert} and scientific text~\citep{beltagy2019scibert}, extending ULMFiT~\citep{howard2018ulmfit}; yet this has not been disentangled for safety-critical generalization. Intermediate-task transfer (STILTs~\citep{phang2018stilts}) and multi-task learning~\citep{caruana1997multitask,liu2019mtdnn} confound domain exposure with supervised signal. Our MLM control confirms pretraining alone closes the OOD gap. Dataset artifact analysis~\citep{gururangan2018artifacts,niven2019probing} motivates the TF-IDF baseline, adversarial NLP~\citep{jin2020textfooler,li2020bertattack,nasr2025attackermovessecond,andriushchenko2024adaptive} informs the perturbation protocol, and probing~\citep{conneau2018probing,hewitt2019structural} underpins CKA analysis.
